\section{Motivation}
Obwohl kinematische Modelle in vielen Bereichen der Robotik eine zentrale Rolle spielen, fällt es Studierenden häufig schwer,
diese abstrakten mathematischen Konzepte vollständig zu durchdringen. Besonders Rotationen im dreidimensionalen Raum – etwa durch Eulerwinkel,
Roll-Pitch-Yaw-Darstellungen, Rotationsmatrizen oder Quaternionen beschrieben – wirken in der Theorie oft wenig anschaulich und sind ohne visuelle
Unterstützung schwer intuitiv erfassbar \cite[94-95]{The_Second_Handbook_of_Research_on_the_Psychology_of_Mathematics_Education}.

In der Hochschullehre entsteht dadurch eine Lücke zwischen theoretischem Wissen und praktischem Verständnis.
Lehrende stehen vor der Herausforderung, komplexe Zusammenhänge nicht nur formal korrekt, sondern auch nachvollziehbar und greifbar zu vermitteln.
Gleichzeitig benötigen Studierende Möglichkeiten, das Gelernte interaktiv zu erproben, um eine tiefere, anwendungsorientierte Intuition für die Kinematik zu entwickeln.

Moderne Lehrmethoden, die auf Visualisierung und Interaktion setzen, können hier entscheidend zur Verbesserung des Lernerfolgs beitragen.
Insbesondere browserbasierte und plattformunabhängige Lösungen, wie sie in JupyterLab realisierbar sind,
eröffnen neue Perspektiven für eine zeitgemäße und praxisnahe Vermittlung technischer Inhalte \cite[239-240,407]{The_Second_Handbook_of_Research_on_the_Psychology_of_Mathematics_Education}.
